\chapter{Die zweiunddreißig Bände}

\section{Übersicht}

\begin{center}
{\small
\renewcommand{\arraystretch}{1.42}
\begin{longtable}{@{}P{1.4cm}P{4.2cm}P{4.2cm}P{1.3cm}@{}}
\toprule
\textbf{Band} & \textbf{Titel} & \textbf{Gebiet} & \textbf{Seiten}\\
\midrule
\endhead
$0$ & Die Kunst des Rechnens & Vorkurs Schulmathematik & $154$\\
$1$ & Die Königin der Mathematik & Zahlentheorie & $181$\\
$2$ & Die Kunst der linearen Struktur & Lineare Algebra & $167$\\
$3$ & Die Kunst des Grenzübergangs & Analysis & $173$\\
$4$ & Die Vermessung des Zufalls & Wahrscheinlichkeitstheorie & $164$\\
$5$ & Die Kunst der Näherung & Numerische Mathematik & $224$\\
$6$ & Die Vermessung des Raumes & Geometrie & $184$\\
$7$ & Die Wege des Denkens & Geschichte der Mathematik & $188$\\
$8$ & Die Kunst der Verformung & Topologie & $178$\\
$9$ & Die Sprache der Felder & Vektor- und Tensoranalysis & $130$\\
$10$ & Die Ordnung des Unendlichen & Logik und Mengenlehre & $124$\\
$11$ & Die Gestalt der Symmetrie & Algebra & $150$\\
$12$ & Die Vollendung der Analysis & Funktionentheorie & $176$\\
$13$ & Das Maß und der Raum & Maßtheorie, Funktionalanalysis & $175$\\
$14$ & Die Gesetze der Veränderung & Differentialgleichungen & $214$\\
$15$ & Die Kunst des Zählens & Kombinatorik & $200$\\
\midrule
$16$ & Punkte und Linien & Graphentheorie & $270$\\
$17$ & Was die Zahlen hergeben & Statistik & $282$\\
$18$ & Wachsen und Vergehen & Biomathematik & $262$\\
$19$ & Das Beste finden & Optimierung & $252$\\
$20$ & Was Nachrichten wiegen & Information, Kryptographie & $242$\\
$21$ & Wenn andere mitrechnen & Spieltheorie & $236$\\
$22$ & Was sich verallgemeinern lässt & Maschinelles Lernen & $204$\\
$23$ & Woran es liegt & Kausale Schlüsse & $158$\\
$24$ & Was zu tun ist & Entscheidungen unter Unsicherheit & $140$\\
$25$ & Wenn der Zufall weitergeht & Stochastische Prozesse & $158$\\
$26$ & Wenn die Daten sich selbst bestimmen & Ökonometrie & $156$\\
$27$ & Was der Rechner vom Zufall weiß & Numerische Simulation & $170$\\
$28$ & Wo die Vorhersage endet & Dynamische Systeme und Chaos & $166$\\
\midrule
$29$ & Die Form der Lösungen & Algebraische Geometrie & $188$\\
$30$ & Was die Symmetrie verrät & Darstellungstheorie & $196$\\
\midrule
$B$ & Die Architektur der Gewissheit & Beweistheorie (Begleitband) & $182$\\
\midrule
\textbf{Gesamt} & & & $\mathbf{6044}$\\
\bottomrule
\end{longtable}}
\end{center}

\section{Die beiden Hälften des Werkes}

\begin{merksatz}
\boxtitel{Ein Unterschied, der sich im Aufbau zeigt}
\begin{center}
{\small
\renewcommand{\arraystretch}{1.55}
\begin{tabular}{@{}P{2.2cm}P{3.0cm}P{3.2cm}@{}}
\toprule
& \textbf{Bände 0 bis 15, 29, 30} & \textbf{Bände 16 bis 28}\\
\midrule
Aufbau & systematisch nach Gebiet & nach Anwendung\\
Reihenfolge & aufeinander aufbauend & je Band eigenständig\\
Beispiele & wechselnd & \textcolor{akzent}{ein durchgehender Leitfall}\\
Ziel & den Stoff verfügbar machen & ihn benutzen können\\
Bezüge & innerhalb des Gebiets & quer durch die erste Hälfte\\
\bottomrule
\end{tabular}}
\end{center}

\begin{center}
\begin{minipage}{0.9\textwidth}
\textbf{Die erste Hälfte fragt: Was ist das und warum stimmt es? Die zweite: Wozu ist es gut und wann versagt es?}
\end{minipage}
\end{center}
\end{merksatz}

\begin{merksatz}
\boxtitel{Der Leitfall}
Jeder der zwölf Anwendungsbände hat ein durchgehendes Beispiel, das vom ersten bis zum letzten Kapitel mitläuft.

\begin{center}
{\small
\renewcommand{\arraystretch}{1.5}
\begin{tabular}{@{}P{1.2cm}P{7.0cm}@{}}
\toprule
$16$ & Das Kontaktnetz einer Station --- von der Gradfolge bis zur Ausbreitung.\\
$17$ & Die MOBIL-Studie zum Hüftgelenksersatz --- von der Beschreibung bis zur Metaanalyse.\\
$18$ & Neustadt und sein Klinikum --- Epidemie, Pharmakokinetik, Tumorwachstum, Resistenz.\\
$19$ & Zwei Eingriffsarten unter knapper OP-Zeit --- von der Ecke bis zum Schattenpreis.\\
$20$ & Die Digitalisierung des Klinikums --- Diagnoseschlüssel, Funkstrecke, Arztbrief.\\
$21$ & Klinikum und Marienhospital --- Standort, Reserveantibiotikum, Dienstplan.\\
$22$ & Die Sepsis-Frühwarnung --- vom Laktatwert bis zum tiefen Netz.\\
$23$ & Die frühe Antibiotikagabe --- eine Tabelle mit falschem Vorzeichen.\\
$24$ & Frau B.\ mit Vorhofflimmern --- Schlaganfall gegen Blutung.\\
$25$ & Die Intensivstation als Prozess --- Ankünfte, Belegung, Wartezeiten.\\
$26$ & Zuzahlung und Arztbesuch --- Experiment, Praxisgebühr, Panel.\\
$27$ & Dieselbe Intensivstation im Rechner --- von der Zufallszahl bis zur Unsicherheit des Ergebnisses.\\
$28$ & Die Herzzelle --- Ruhe, Takt, Flimmern, und wo die Vorhersage endet.\\
\bottomrule
\end{tabular}}
\end{center}

\begin{center}
\begin{minipage}{0.88\textwidth}
\kursiv{Die Wahl des Gegenstands ist nicht zufällig: \textbf{Ein Krankenhaus ist ein Ort, an dem Zahlen Folgen haben} --- und an dem sich gut zeigen lässt, warum eine falsche Modellannahme nicht nur ein Rechenfehler ist. Wer aus einem anderen Fach kommt, kann die Leitfälle ohne Vorkenntnisse lesen; sie sind durchgehend erklärt.}
\end{minipage}
\end{center}
\end{merksatz}

\begin{merksatz}
\boxtitel{Wie die Anwendungsbände zusammenhängen}
\begin{center}
\begin{tikzpicture}[scale=1.0]
\node[knotenb] (b16) at (-3.4,1.5) {16};
\node[knotenb] (b17) at (-1.1,1.5) {17};
\node[knotenb] (b18) at (1.2,1.5) {18};
\node[knotenb] (b19) at (-2.2,0) {19};
\node[knotenb] (b20) at (0.1,0) {20};
\node[knotenb] (b21) at (2.4,0) {21};
\node[knotenb] (b22) at (0.1,-1.5) {22};
\node[knotenb] (b23) at (2.6,-1.5) {23};
\node[knotenb] (b24) at (5.0,-1.5) {24};
\node[knotenr] (b25) at (-2.2,-1.5) {25};
\node[knotenb] (b26) at (1.2,-3.0) {26};
\node[knotenb] (b27) at (-3.4,-3.0) {27};
\node[knotenb] (b28) at (3.8,-3.0) {28};
\draw[kanted, -{Stealth[length=4pt]}] (b17)--(b18);
\draw[kanted, -{Stealth[length=4pt]}] (b17)--(b20);
\draw[kanted, -{Stealth[length=4pt]}] (b16)--(b19);
\draw[kanted, -{Stealth[length=4pt]}] (b19)--(b20);
\draw[kanted, -{Stealth[length=4pt]}] (b19)--(b21);
\draw[kanted, -{Stealth[length=4pt]}] (b20)--(b21);
\draw[kanted, -{Stealth[length=4pt]}] (b18) to[bend left=22] (b21);
\draw[kanted, -{Stealth[length=4pt]}] (b17) to[bend right=35] (b22);
\draw[kanted, -{Stealth[length=4pt]}] (b19) to[bend right=15] (b22);
\draw[kanted, -{Stealth[length=4pt]}] (b20)--(b22);
\draw[kanted, -{Stealth[length=4pt]}] (b22)--(b23);
\draw[kanted, -{Stealth[length=4pt]}] (b23)--(b24);
\draw[kanted, -{Stealth[length=4pt]}] (b19) to[bend left=12] (b24);
\draw[kanted, -{Stealth[length=4pt]}] (b21) to[bend left=25] (b24);
\draw[kanted, -{Stealth[length=4pt]}] (b25) to[bend left=10] (b22);
\draw[kanted, -{Stealth[length=4pt]}] (b25) to[bend right=25] (b18);
\draw[kanted, -{Stealth[length=4pt]}] (b25)--(b23);
\draw[kanted, -{Stealth[length=4pt]}] (b17) to[bend right=50] (b26);
\draw[kanted, -{Stealth[length=4pt]}] (b22)--(b26);
\draw[kanted, -{Stealth[length=4pt]}] (b23)--(b26);
\draw[kanted, -{Stealth[length=4pt]}] (b25)--(b26);
\draw[kanted, -{Stealth[length=4pt]}] (b25)--(b27);
\draw[kanted, -{Stealth[length=4pt]}] (b17) to[bend right=40] (b27);
\draw[kanted, -{Stealth[length=4pt]}] (b25) to[bend right=12] (b28);
\draw[kanted, -{Stealth[length=4pt]}] (b18) to[bend left=40] (b28);
\draw[kanted, -{Stealth[length=4pt]}] (b16) to[bend left=20] (b23);
\draw[kanted, -{Stealth[length=4pt]}] (b17) to[bend right=8] (b23);

\node[font=\scriptsize\itshape, color=grau] at (-1.6,-3.6) {ein Pfeil: der eine Band benutzt den anderen};
\end{tikzpicture}
\end{center}

\begin{center}
{\small
\renewcommand{\arraystretch}{1.5}
\begin{tabular}{@{}P{2.0cm}P{6.4cm}@{}}
\toprule
$17 \to 18$ & Schätzen und Testen für Modellanpassung und Studienauswertung.\\
$17 \to 20$ & Die relative Entropie erscheint dort als Werkzeug der Modellwahl.\\
$16 \to 19$ & Flüsse und Zuordnungen als ganzzahlige lineare Programme.\\
$19 \to 20$ & Dualität und Konvexität hinter der Kanalkapazität.\\
$19 \to 21$ & Der Minimax-Satz ist der Dualitätssatz.\\
$20 \to 21$ & Robuste Optimierung als Spiel gegen die Natur.\\
$18 \to 21$ & Resistenz als evolutionäres Spiel --- die Rechnung, die Band 18 offenließ.\\
$17,\,19,\,20 \to 22$ & Schätzen, Minimieren und die kürzeste Beschreibung treffen sich im Lernen.\\
$16,\,17,\,22 \to 23$ & Graphen tragen die Annahmen, das Lernen die Schätzer.\\
$19,\,21,\,23 \to 24$ & Wirkung, Kosten und Gegenspieler treffen sich in der Wahl.\\
$25 \to 18,\,22,\,23$ & Der einzige Band der zweiten Hälfte, der Grundlagen nachliefert statt anzuwenden.\\
$17,\,22,\,23,$ $25 \to 26$ & Der Band, der die meisten Vorgänger zugleich braucht --- und sie unter erschwerten Bedingungen zusammenführt.\\
\bottomrule
\end{tabular}}
\end{center}

\begin{center}
\begin{minipage}{0.88\textwidth}
\kursiv{\textbf{Die letzte Zeile ist die engste Verbindung des ganzen Werkes.} Band 18 behandelt in Kapitel~14 die adaptive Tumortherapie und bemerkt, das Wettrüsten zwischen Therapie und Resistenz sei im Kern ein Spiel --- ohne die Rechnung zu führen. Band 21 führt sie in Kapitel~10 mit denselben Zahlen aus: Startanteil $10^{-6}$, Eignungsdifferenzen $+0{,}04$ und $-0{,}01$ je Tag, Ergebnis $T_{\text{aus}} \ge 4\,T_{\text{an}}$.}
\end{minipage}
\end{center}
\end{merksatz}

\section{Was in welchem Band steht}

\begin{merksatz}
\boxtitel{Band 0 --- Die Kunst des Rechnens}
Die gesamte Schulmathematik, jede Regel begründet. Zahlbereiche und der Irrationalitätsbeweis für $\sqrt2$; Bruchrechnung mit drei Begründungen für den Kehrwert; Potenzen, Wurzeln, Logarithmen als eine Familie; Terme, Gleichungen, Ungleichungen mit Beträgen; der Funktionsbegriff und die acht Grundtypen; Geometrie und Trigonometrie; Folgen, Ableitung, Integral; Beweisen lernen.

\kursiv{Besonderheit: ein Fehlerkatalog mit dreiunddreißig Einträgen und ein Schlusskapitel über die Lernforschung.}
\end{merksatz}

\begin{merksatz}
\boxtitel{Band 1 --- Die Königin der Mathematik}
Zahlentheorie von der Teilbarkeit bis zu elliptischen Kurven. Euklidischer Algorithmus, Fundamentalsatz der Arithmetik, modulares Rechnen, chinesischer Restsatz, die Sätze von Fermat, Euler und Wilson, quadratische Reziprozität, Primzahlverteilung, diophantische Gleichungen, Kettenbrüche, RSA und Primzahltests.

\kursiv{Der Titel stammt von Gauß, der die Zahlentheorie die Königin der Mathematik nannte.}
\end{merksatz}

\begin{merksatz}
\boxtitel{Band 2 --- Die Kunst der linearen Struktur}
Lineare Algebra vom Gleichungssystem bis zur Singulärwertzerlegung. Der Aufbau geht bewusst vom Konkreten zum Abstrakten: erst Gauß-Algorithmus und Matrizen, dann der Vektorraumbegriff, dann lineare Abbildungen, Determinante, Eigenwerte, Jordanform, Skalarprodukte, Spektralsatz.

\kursiv{Die Singulärwertzerlegung ist der Grund, warum dieser Band für alle Anwendungsgebiete zentral ist --- von der Bildkompression bis zum maschinellen Lernen.}
\end{merksatz}

\begin{merksatz}
\boxtitel{Band 3 --- Die Kunst des Grenzübergangs}
Analysis in der strengen Fassung. Konstruktion der reellen Zahlen, Folgen und Reihen, Stetigkeit, die elementaren Funktionen, Differentialrechnung mit Mittelwertsätzen und Taylor, Riemann-Integral und Hauptsatz, gleichmäßige Konvergenz, Differentialrechnung im $\R^{n}$.

\kursiv{Der Unterschied zu Band 0: Dort wird gerechnet, hier wird bewiesen. Die $\varepsilon$-$\delta$-Definition ist der Angelpunkt.}
\end{merksatz}

\begin{merksatz}
\boxtitel{Band 4 --- Die Vermessung des Zufalls}
Wahrscheinlichkeitstheorie von null an, einschließlich der Schulstochastik. Kombinatorik, bedingte Wahrscheinlichkeit und das Ziegenproblem, Zufallsvariablen, die diskreten und stetigen Verteilungen, Erwartungswert und Varianz, die Gesetze der großen Zahlen, der zentrale Grenzwertsatz, Markow-Ketten, Martingale, Schätzen und Testen.

\kursiv{Wer die Schulstochastik wiederholen will, findet sie hier --- deshalb fehlt sie in Band 0.}
\end{merksatz}

\begin{merksatz}
\boxtitel{Band 5 --- Die Kunst der Näherung}
Numerische Mathematik. Gleitkommazahlen und ihre Tücken, Kondition und Stabilität, LR- und QR-Zerlegung, iterative Verfahren, Eigenwertberechnung, Nullstellenverfahren, Interpolation und Splines, Quadratur, Anfangs- und Randwertprobleme, Optimierung, Monte-Carlo.

\kursiv{Der Band beginnt mit einer Warnung: $0{,}1+0{,}2 \ne 0{,}3$ im Rechner --- und erklärt, warum das so ist und was daraus folgt.}
\end{merksatz}

\begin{merksatz}
\boxtitel{Band 6 --- Die Vermessung des Raumes}
Geometrie in allen ihren Gestalten. Euklid und die Axiome, Dreieck und Kegelschnitte, analytische Geometrie, Isometrien und Symmetriegruppen, projektive Geometrie, das Parallelenproblem, hyperbolische und sphärische Geometrie, Kurven- und Flächentheorie bis Gauß-Bonnet, diskrete Geometrie und Packungen.

\kursiv{Das Parallelenproblem ist die Achse des Bandes --- zweitausend Jahre Bemühen um einen Beweis, der nicht zu führen war.}
\end{merksatz}

\begin{merksatz}
\boxtitel{Band 7 --- Die Wege des Denkens}
Die Geschichte der Mathematik im Zusammenhang. Von den Zahlzeichen über China, Indien und Mesoamerika, das griechische Wunder, die islamische Welt, Renaissance, die Revolution des 17.\ Jahrhunderts, die Rigorisierung des 19., die Grundlagenkrise des 20.\ --- bis zu Vertreibung, Krieg und der Mathematik von heute.

\kursiv{Der Band ist ohne Vorkenntnisse lesbar und eignet sich als Einstieg für alle, die wissen wollen, worum es überhaupt geht.}
\end{merksatz}

\begin{merksatz}
\boxtitel{Band 8 --- Die Kunst der Verformung}
Topologie vom metrischen Raum bis zur Homologie. Offene Mengen, Stetigkeit ohne $\varepsilon$, Kompaktheit und Zusammenhang, Trennungsaxiome, Quotientenräume, Fundamentalgruppe und Überlagerungen, Homologie, Fixpunktsätze, Mannigfaltigkeiten und die Klassifikation der Flächen.

\kursiv{Der Fixpunktsatz von Brouwer taucht in vier weiteren Bänden wieder auf --- zuletzt in Band 21 als Grundlage des Existenzsatzes für Gleichgewichte.}
\end{merksatz}

\begin{merksatz}
\boxtitel{Band 9 --- Die Sprache der Felder}
Vektor- und Tensoranalysis. Skalar- und Vektorfelder, Gradient, Divergenz, Rotation, Kurven- und Flächenintegrale, die Sätze von Green, Gauß und Stokes, krummlinige Koordinaten, Tensoren, Differentialformen und der allgemeine Satz von Stokes.

\kursiv{Die Pointe des Bandes: Die drei Integralsätze sind ein einziger Satz, in drei Gestalten.}
\end{merksatz}

\begin{merksatz}
\boxtitel{Band 10 --- Die Ordnung des Unendlichen}
Logik und Mengenlehre. Aussagen- und Prädikatenlogik, Beweiskalküle, Vollständigkeits- und Kompaktheitssatz, ZFC mit allen Axiomen, Ordinal- und Kardinalzahlen, Auswahlaxiom und seine Äquivalente, Berechenbarkeit, die Unvollständigkeitssätze, Unabhängigkeit und die Kontinuumshypothese.

\kursiv{Überschneidet sich mit dem Begleitband B: Band 10 liefert die Technik, der Begleitband die Einordnung.}
\end{merksatz}

\begin{merksatz}
\boxtitel{Band 11 --- Die Gestalt der Symmetrie}
Algebra von der Gruppe bis zur Galoistheorie. Gruppen, Untergruppen, Satz von Lagrange, Normalteiler und Faktorgruppen, Gruppenoperationen und Sylowsätze, Ringe und Ideale, Polynomringe, Körpererweiterungen, Konstruierbarkeit mit Zirkel und Lineal, endliche Körper, die Galoiskorrespondenz und die Unlösbarkeit der Quintik.

\kursiv{Die drei klassischen Unmöglichkeiten --- Würfelverdopplung, Winkeldreiteilung, Quadratur des Kreises --- werden vollständig bewiesen.}
\end{merksatz}

\begin{merksatz}
\boxtitel{Band 12 --- Die Vollendung der Analysis}
Funktionentheorie. Komplexe Zahlen, Holomorphie und die Cauchy-Riemann-Gleichungen, der Cauchysche Integralsatz, Potenz- und Laurentreihen, der Residuensatz mit Anwendungen auf reelle Integrale, konforme Abbildungen, der Riemannsche Abbildungssatz, analytische Fortsetzung, Gamma- und Zetafunktion.

\kursiv{Der Band heißt so, weil in ihm zusammenläuft, was in den Bänden 2, 3 und 9 getrennt entwickelt wurde.}
\end{merksatz}

\begin{merksatz}
\boxtitel{Band 13 --- Das Maß und der Raum}
Maßtheorie und Funktionalanalysis. Die drei Defekte des Riemann-Integrals, $\sigma$-Algebren und Maße, das Lebesgue-Maß und die Menge von Vitali, messbare Funktionen, die Konvergenzsätze, Fubini, die $L^{p}$-Räume, Banach- und Hilberträume, Hahn-Banach und die drei Prinzipien, kompakte Operatoren, der Spektralsatz.

\kursiv{Das Schlusskapitel zeigt, warum die Quantenmechanik ohne diesen Band nicht formulierbar wäre.}
\end{merksatz}

\begin{merksatz}
\boxtitel{Band 14 --- Die Gesetze der Veränderung}
Differentialgleichungen. Richtungsfeld und elementare Lösungsverfahren, Existenz und Eindeutigkeit nach Picard-Lindelöf, lineare Systeme und das Matrixexponential, Schwingungen und Resonanz, Sturm-Liouville, Phasenporträt und Stabilität, Bifurkationen und Chaos, die drei klassischen partiellen Gleichungen, schwache Lösungen und finite Elemente.

\kursiv{Der Band ist die Brücke zwischen Analysis und Anwendung --- und die unmittelbare Grundlage von Band 18.}
\end{merksatz}

\begin{merksatz}
\boxtitel{Band 15 --- Die Kunst des Zählens}
Kombinatorik. Die Grundregeln und das Schubfachprinzip, die vier Grundaufgaben, das Pascalsche Dreieck, Inklusion und Exklusion, erzeugende Funktionen, Catalan- und Stirlingzahlen, Partitionen, Burnside und Pólya, Graphen, Ramseytheorie, die probabilistische Methode, extremale Mengenlehre.

\kursiv{Der Graphenteil ist hier ein Kapitel; Band 16 macht ein eigenes Buch daraus.}
\end{merksatz}

\begin{merksatz}
\boxtitel{Band 16 --- Punkte und Linien}
Graphentheorie als eigenständiger Band. Vom Graphbegriff über Wege und Zusammenhang, Bäume und Spannbäume, Eulerkreise und Hamiltonkreise, Planarität und der Vierfarbensatz, Färbungen, Matchings und der Satz von Hall, Flüsse und Schnitte, Zufallsgraphen, Netzwerke und Zentralitätsmaße.

\kursiv{Leitfall ist das Kontaktnetz einer Station --- vom Handschlaglemma bis zur Ausbreitungsschwelle. Der erste Band der zweiten Werkhälfte.}
\end{merksatz}

\begin{merksatz}
\boxtitel{Band 17 --- Was die Zahlen hergeben}
Statistik von der Beschreibung bis zur Metaanalyse. Lage- und Streumaße, Schätzen, Konfidenzintervalle, Tests und ihre Fehler, Fallzahlplanung, lineare und logistische Regression, Überlebenszeitanalyse, Bayessche Verfahren, Studiendesign, multiples Testen, Metaanalyse.

\kursiv{Leitfall ist die MOBIL-Studie zum Hüftgelenksersatz. Der Band nimmt die Replikationskrise zum Anlass, jede Methode auch daran zu messen, wie leicht sie sich missbrauchen lässt.}
\end{merksatz}

\begin{merksatz}
\boxtitel{Band 18 --- Wachsen und Vergehen}
Biomathematik. Exponentielles und logistisches Wachstum, Chaos in der Rekursion, Lotka-Volterra und Wettbewerbsmodelle, Stabilität, SIR-Modelle und Impfschwellen, Ausbreitung auf Netzwerken, Pharmakokinetik, Enzymkinetik, erregbare Systeme, Tumorwachstum nach Gompertz, Resistenz und adaptive Therapie.

\kursiv{Leitfall ist Neustadt mit seinen $100\,000$ Einwohnern und seinem Klinikum. Der Band benutzt die Bände 14 und 17 unmittelbar.}
\end{merksatz}

\begin{merksatz}
\boxtitel{Band 19 --- Das Beste finden}
Optimierung. Lineare Programme und das Simplexverfahren, Dualität und Schattenpreise, totale Unimodularität, Konvexität, Lagrange und Karush-Kuhn-Tucker, Gradienten- und Newtonverfahren, ganzzahlige Optimierung, Verzweigen und Beschränken, Komplexität, Optimierung unter Unsicherheit, mehrere Ziele.

\kursiv{Leitfall ist die Planung zweier Eingriffsarten unter knapper OP-Zeit. Kapitel~4 enthält die Bemerkung, die Band 21 aufgreift: dass von Neumann in Dantzigs Problem sein eigenes wiedererkannte.}
\end{merksatz}

\begin{merksatz}
\boxtitel{Band 20 --- Was Nachrichten wiegen}
Informationstheorie und Kryptographie in einem Band. Informationsgehalt und Entropie, Quellencodierung und Huffman, Transinformation und Kanalkapazität, das Kanalcodierungstheorem, fehlerkorrigierende Codes, Kolmogorow-Komplexität, Information in der Statistik, perfekte Sicherheit, Blockchiffren, Hashfunktionen, Diffie-Hellman, RSA, elliptische Kurven.

\kursiv{Die Klammer des Bandes: Dieselbe Größe misst, was ankommen soll, und was nicht ankommen darf.}
\end{merksatz}

\begin{merksatz}
\boxtitel{Band 21 --- Wenn andere mitrechnen}
Spieltheorie. Nullsummenspiele und der Minimax-Satz, Spiele als lineare Programme, das Nash-Gleichgewicht, die klassischen Dilemmata, Existenz und Berechenbarkeit, extensive Form und Glaubwürdigkeit, wiederholte Spiele, evolutionär stabile Strategien, Replikatordynamik, Kollektivgüter, Abstimmungsverfahren, Auktionen, Verhandlung und Koalitionen.

\kursiv{Leitfall sind Klinikum und Marienhospital in Neustadt. Der Band schließt die Rechnung ab, die Band 18 offenließ.}
\end{merksatz}

\begin{merksatz}
\boxtitel{Band 22 --- Was sich verallgemeinern lässt}
Maschinelles Lernen. Die Lernaufgabe und die Bayes-Regel, Verzerrung und Varianz, Kreuzvalidierung und Datenlecks, Regularisierung mit Kamm und Lasso, logistische Regression und Perzeptron, Stützvektoren und Kerne, PAC-Lernen und VC-Dimension, Verallgemeinerungsschranken, Bäume und Wälder, nächste Nachbarn, neuronale Netze, Gruppenbildung, Hauptkomponenten, tiefe Netze und der doppelte Abstieg.

\kursiv{Leitfall ist die Sepsis-Frühwarnung --- und als Maßstab die Nachprüfung eines kommerziellen Warnsystems von 2021, die seinen Fehlschlag offenlegte. Der Band benutzt die Bände 17, 19 und 20 durchgehend.}
\end{merksatz}

\begin{merksatz}
\boxtitel{Band 23 --- Woran es liegt}
Kausale Schlüsse. Zusammenhang und Ursache, potentielle Ergebnisse, Randomisierung, kausale Graphen und d-Trennung, zulässige Anpassungsmengen, Hintertür- und Vordertürformel, Behandlungsneigung, Gewichtung und doppelte Absicherung, Instrumentvariablen, Regressionssprung und doppelte Differenz, Zielstudienemulation, g-Formel, Mediation, heterogene Effekte.

\kursiv{Leitfall ist die frühe Antibiotikagabe bei Sepsis --- eine Tabelle, die das Gegenteil dessen zeigt, was die Leitlinien sagen. Der Band benutzt die Bände 16, 17 und 22 und schließt die zweite Werkhälfte ab.}
\end{merksatz}

\begin{merksatz}
\boxtitel{Band 24 --- Was zu tun ist}
Entscheidungen unter Unsicherheit. Erwartungsnutzen und Behandlungsschwelle, Nutzenaxiome und Risikoscheu, Wert der Information, Test- und Behandlungsschwellen, Entscheidungsbäume und Empfindlichkeitsanalyse, Markow-Modelle und QALY, Rückwärtsrechnung, Markow-Entscheidungsprozesse, Erkunden gegen Ausnutzen, Kosten-Wirksamkeit, Ambiguität, tatsächliches Entscheidungsverhalten, Gruppenentscheidungen, Priorisierung.

\kursiv{Leitfall ist Frau B.\ mit Vorhofflimmern: Schlaganfall gegen Blutung. Der Band benutzt die Bände 19, 21 und 23 und schließt an die Frage an, die der kausale Schluss offenlässt.}
\end{merksatz}

\begin{merksatz}
\boxtitel{Band 25 --- Wenn der Zufall weitergeht}
Stochastische Prozesse. Irrfahrten und das Arkussinus-Gesetz, Markow-Ketten und stationäre Verteilungen, Poisson-Prozesse, Warteschlangen, Verzweigungsprozesse, bedingte Erwartung und Filtration, Martingale und optionales Stoppen, Konvergenzsätze, Brownsche Bewegung, Itō-Integral und Itō-Formel, Optionsbewertung, Prozesse in Klinik und Natur.

\kursiv{Der einzige Band der zweiten Werkhälfte, der eine Grundlagenlücke schließt: Sechs frühere Bände hatten den Prozessbegriff gebraucht und jeweils nur so weit entwickelt, wie es gerade nötig war.}
\end{merksatz}

\begin{merksatz}
\boxtitel{Band 26 --- Wenn die Daten sich selbst bestimmen}
Ökonometrie. Endogenität, Identifikation und Abzählbedingung, Instrumentvariablen und schwache Instrumente, verallgemeinerte Momentenmethode, Logit, Probit und Tobit, Auswahlmodelle, feste Effekte, dynamische Panels, doppelte Differenz und gestaffelte Einführung, Stationarität, Scheinregression und Einheitswurzeln, Kointegration und Fehlerkorrektur, Vektorautoregression, Lucas-Kritik, Vorhersage gegen Politik.

\kursiv{Leitfall ist die Wirkung von Zuzahlungen auf Arztbesuche, gemessen am Krankenversicherungsexperiment und an der Praxisgebühr. Der Band verbindet die Bände 17, 22, 23 und 25 --- kausale Fragen, Beobachtungsdaten, rückgekoppelte Systeme und nichtstationäre Zeitreihen zugleich.}
\end{merksatz}

\begin{merksatz}
\boxtitel{Band 27 --- Was der Rechner vom Zufall weiß}
Numerische Simulation. Pseudozufallszahlen und ihre Prüfung, Inversion, Verwerfung und Box--Muller, Monte-Carlo-Integration und der Fluch der Dimension, Standardfehler bei Abhängigkeit, Varianzreduktion (antithetisch, Kontrollvariable, Schichtung), Wichtigkeitsstichprobe, Metropolis und Gibbs, Konvergenzdiagnostik, Hamilton-Monte-Carlo, Pfade stochastischer Differentialgleichungen, ereignisgesteuerte Simulation, seltene Ereignisse, Teilchenfilter, Reproduzierbarkeit, Unsicherheit des Ergebnisses.

\kursiv{Leitfall ist die Intensivstation aus Band 25, diesmal als Programm: Jede Tabelle des Bandes entsteht aus einem Skript, das beiliegt. Der Band benutzt die Bände 4, 5, 17 und 25 und ist der Band, der die Zahlen der übrigen Anwendungsbände nachrechenbar macht.}
\end{merksatz}

\begin{merksatz}
\boxtitel{Band 28 --- Wo die Vorhersage endet}
Dynamische Systeme und Chaos. Fluss, Fixpunkte und Stabilität, Grenzzyklen und Poincaré--Bendixson, Bifurkationen (Sattel-Knoten, Heugabel, Hopf), erzwungene Schwingungen und Synchronisation, Periodenverdopplung und Feigenbaum, Lorenz-System, Lyapunov-Exponenten und Vorhersagehorizont, Einbettung nach Takens, fraktale Dimension, Surrogatdaten, Reaktions-Diffusions-Muster und Spiralwellen, gekoppelte Oszillatoren und Kuramoto.

\kursiv{Leitfall ist die Herzzelle: vom Ruhepotential über den Schrittmacher bis zum Kammerflimmern als Spiralwelle. Der Band benutzt die Bände 14, 18 und 25 und zeigt, warum deterministisch nicht vorhersagbar heißt.}
\end{merksatz}

\begin{merksatz}
\boxtitel{Band 29 --- Die Form der Lösungen}
Algebraische Geometrie. Varietäten und Ideale, Hilberts Basissatz und Nullstellensatz, Monomordnungen und Division, Gröbner-Basen und Buchbergers Algorithmus, Elimination und Resultanten, Hilbert-Funktion, Dimension und Grad, Tangentialraum und Glattheit, projektiver Raum und Bézout, elliptische Kurven und ihr Gruppengesetz, rationale Parametrisierung, Geschlecht, algebraische Statistik, Kinematik und Kryptographie.

\kursiv{Leitfall ist die Identifizierbarkeit eines Zwei-Kompartiment-Modells: Welche Parameter bestimmen die Messkurve? Ein Grundlagenband, geschrieben nach der zweiten Werkhälfte; er benutzt die Bände 2, 11 und 12 und beantwortet eine Frage, die Band 18 bei der Pharmakokinetik offen lassen musste.}
\end{merksatz}

\begin{merksatz}
\boxtitel{Band 30 --- Was die Symmetrie verrät}
Darstellungstheorie. Punktgruppen und Symmetrieoperationen, Darstellungen und Unitarität, Maschke und Schur, Charaktere und Orthogonalität, Charaktertafeln, Normalschwingungen und Auswahlregeln, Einschränkung und Induktion, Frobenius-Reziprozität, symmetrische Gruppe und Young-Diagramme, Tensorprodukte und Frobenius--Schur-Indikator, $SO(3)$ und Kugelflächenfunktionen, $SU(2)$ und Spin, Fourier auf Gruppen, Peter--Weyl, Schur--Weyl, Hückel-Orbitale des Benzols, Kristalle, Kristallfeld am Hämoglobin, Quarks.

\kursiv{Leitfall ist das Wassermolekül: vier Symmetrieoperationen, neun Bewegungen, drei Schwingungen, und was man davon im Spektrum sieht. Der Band benutzt die Bände 2, 11 und 13 und erklärt, warum Orbitale zu $1, 3, 5, 7$ kommen.}
\end{merksatz}

\begin{merksatz}
\boxtitel{Band B --- Die Architektur der Gewissheit}
Der Begleitband über den Beweisbegriff selbst. Von der Welt vor Griechenland über Euklid und Archimedes, die Krise der Infinitesimalien, die Rigorisierung, Cantor und die Grundlagenkrise; Aussagen- und Prädikatenlogik, ZFC; die Beweisformen und sechs Beweise unter dem Mikroskop; Berechenbarkeit, Gödel, Unabhängigkeit, große Kardinalzahlen, Beweiskomplexität; konstruktive Mathematik, Typentheorie, Kategorien; Computerbeweise und Beweisassistenten.

\kursiv{Der erste geschriebene Band und der einzige ohne Nummer.}
\end{merksatz}

\section{Lesewege}

\begin{merksatz}
\boxtitel{Sieben Wege durch das Werk}
\begin{center}
{\small
\renewcommand{\arraystretch}{1.55}
\begin{tabular}{@{}P{3.2cm}P{5.4cm}@{}}
\toprule
\textbf{Von vorn} & $0 \to 3 \to 2 \to 4 \to 11$ und weiter. Der vollständige Weg; Band 0 zuerst, auch wenn man glaubt, ihn nicht zu brauchen.\\
\textbf{Studienbegleitend} & Zum jeweiligen Semester den passenden Band: erstes Semester $0$, $2$, $3$; zweites $3$, $4$, $11$; drittes $12$, $13$, $14$; viertes $5$, $15$.\\
\textbf{Nur verstehen} & $7 \to B \to 10$. Ohne Rechnen, aber mit Einordnung. Für alle, die wissen wollen, worum es geht.\\
\textbf{Anwendung} & $17 \to 19 \to 21$ oder $17 \to 18$. Die zweite Hälfte lässt sich weitgehend für sich lesen; nötige Grundlagen werden jeweils benannt.\\
\textbf{Medizin und Biologie} & $17 \to 18 \to 22 \to 23 \to 24$, ergänzt um $16$ für Netzwerke, $21$ für Anreize und $4$ für die Wahrscheinlichkeitsrechnung.\\
\textbf{Wirtschaft} & $2 \to 3 \to 4 \to 17 \to 19 \to 21 \to 25 \to 26$, ergänzt um $23$ für die kausalen Begriffe und $24$ für die Entscheidungen.\\
\textbf{Natur\-wissenschaften} & $3 \to 2 \to 14 \to 28$ für die Dynamik, $11 \to 30$ für die Symmetrie, $4 \to 25 \to 27$ für die Simulation; $29$ für alle, die Gleichungssysteme lösen müssen.\\
\bottomrule
\end{tabular}}
\end{center}

\begin{center}
\begin{minipage}{0.88\textwidth}
\kursiv{\textbf{Die Wege vier bis sieben sind erst mit der zweiten Werkhälfte und den Nachtragsbänden 27 bis 30 möglich geworden.} Sie sind erst mit der zweiten Werkhälfte möglich geworden und setzen weniger voraus, als man vermutet: Wer Band 0 beherrscht und die Grundzüge der Wahrscheinlichkeitsrechnung kennt, kann Band 17 lesen --- und von dort aus weiter.}
\end{minipage}
\end{center}
\end{merksatz}

\begin{merksatz}
\boxtitel{Was welcher Band voraussetzt}
\begin{center}
{\small
\renewcommand{\arraystretch}{1.44}
\begin{tabular}{@{}P{1.4cm}P{6.8cm}@{}}
\toprule
$16$ & Band 0; Band 15 nützlich, nicht nötig.\\
$17$ & Band 0 und die Grundzüge von Band 4.\\
$18$ & Band 14 für die Differentialgleichungen, Band 17 für die Auswertung.\\
$19$ & Band 2 für Matrizen, Band 3 für Ableitungen.\\
$20$ & Band 4 für Wahrscheinlichkeiten, Band 19 für die Dualität, Band 1 für die Zahlentheorie ab Kapitel~12.\\
$21$ & Band 19 durchgehend, Band 12 für Band 21 nicht nötig.\\
$22$ & Band 17 für Schätzen und Testen, Band 19 für Konvexität, Band 2 für Matrizen; Band 20 ab Kapitel~9.\\
$23$ & Band 4 für bedingte Wahrscheinlichkeit, Band 16 für Graphen, Band 17 für Schätzen; Teil I und II ohne weitere Voraussetzungen.\\
$24$ & Band 4 für Erwartungswert und Bayes; Band 19 erst ab Kapitel~7, Band 21 nur für Kapitel~13.\\
$25$ & Band 4 durchgehend, Band 2 für Kapitel~2 und 3; Band 13 erst ab Kapitel~7.\\
$26$ & Band 17 durchgehend, Band 23 für die kausalen Begriffe; Band 25 nur für Teil~IV.\\
$27$ & Band 4 durchgehend, Band 25 für die Ketten und Prozesse; Band 5 für Kapitel~3 und 9, Band 17 für Kapitel~4.\\
$28$ & Band 14 für Differentialgleichungen, Band 2 für Eigenwerte; Band 25 nur für Kapitel~12, Band 18 für den Leitfall nützlich.\\
$29$ & Band 11 für Ringe und Ideale, Band 2 für lineare Algebra; Band 12 für Kapitel~9 und 12, Band 1 für Kapitel~14.\\
$30$ & Band 11 für Gruppen, Band 2 für Matrizen und Eigenwerte; Band 13 für Kapitel~10 und 12, Band 25 nur für Verweise.\\
\bottomrule
\end{tabular}}
\end{center}
\end{merksatz}

\section{Der gemeinsame Aufbau}

\begin{merksatz}
\boxtitel{Jeder Band folgt demselben Muster}
\begin{center}
{\small
\renewcommand{\arraystretch}{1.5}
\begin{tabular}{@{}P{2.6cm}P{6.0cm}@{}}
\toprule
\textbf{Vorwort} & Was der Band will, wen er anspricht, was er voraussetzt.\\
\textbf{Teile} & Sechs bis acht, jeder mit eigenem Schwerpunkt.\\
\textbf{Kapitel} & Fünfzehn bis zwanzig, jedes mit demselben inneren Aufbau.\\
\textbf{Aktueller Stand} & Das vorletzte Kapitel: was offen ist.\\
\textbf{Anhänge} & Formelsammlung, Zeittafel, Glossar, kommentierte Quellen.\\
\bottomrule
\end{tabular}}
\end{center}

\schritt{Der Aufbau innerhalb eines Kapitels}
\begin{center}
{\small
\renewcommand{\arraystretch}{1.5}
\begin{tabular}{@{}P{3.0cm}P{5.6cm}@{}}
\toprule
\textbf{Der Leitfall} & Eine Frage, an der sich das Kapitel entwickelt.\\
\textbf{Definitionen} & Zeichen für Zeichen erklärt, mit der Frage, warum gerade so.\\
\textbf{Sätze} & Mit vollständigem Beweis oder mit benannter Auslassung.\\
\textbf{Der Bruchfall} & Wo die Sache nicht mehr trägt.\\
\textbf{Übungsaufgaben} & Drei, vollständig durchgerechnet, mit Probe.\\
\textbf{Geschichte} & Woher es kommt und wer sich geirrt hat.\\
\textbf{Forschung} & Woran heute gearbeitet wird.\\
\bottomrule
\end{tabular}}
\end{center}

\begin{center}
\begin{minipage}{0.88\textwidth}
\kursiv{\textbf{Der Abschnitt \enquote{Der Bruchfall} ist die auffälligste Eigenheit der zweiten Werkhälfte.} Er steht in jedem Kapitel und benennt, was das gerade Entwickelte \kursiv{nicht} leistet --- welche Voraussetzung selten erfüllt ist, welcher Schluss nicht gilt, welcher Fehler naheliegt. In den Bänden 0 bis 15 übernehmen diese Aufgabe die Warnkästen.}
\end{minipage}
\end{center}
\end{merksatz}

\begin{merksatz}
\boxtitel{Die Prüfschwelle}
Jeder Band wurde vor der Freigabe maschinell geprüft:

\begin{center}
{\small
\renewcommand{\arraystretch}{1.48}
\begin{tabular}{@{}P{3.2cm}P{2.0cm}P{2.6cm}@{}}
\toprule
\textbf{Kriterium} & \textbf{Schwelle} & \\
\midrule
Satzfehler & $0$ & \\
Undefinierte Verweise & $0$ & \\
Überlaufende Zeilen & $0$ & über $2$ Punkt\\
Kommentierte Quellen & $28$ & je Band\\
\bottomrule
\end{tabular}}
\end{center}

\begin{center}
\begin{minipage}{0.88\textwidth}
\kursiv{Hinzu kommt eine Prüfung, die sich nicht automatisieren lässt: \textbf{Jede Zahl in jeder Aufgabe wurde nachgerechnet}, und zwar getrennt vom Text. Mehrfach hat das zu Berichtigungen geführt, die im laufenden Text stehengeblieben wären.}
\end{minipage}
\end{center}
\end{merksatz}
